\begin{theorem}[Left adjoints preserve colimits; right adjoints preserve limits]
    \label{theorem:right_left_adjoints_commute_with_limits_colimits}
    \TODO{AI generated, verify}
    Let $\mathcal{C}$ and $\mathcal{D}$ be \CrefAndHyperrefIfExist{definition:category}{categories}. Let $F : \mathcal{C} \to \mathcal{D}$ and $G : \mathcal{D} \to \mathcal{C}$ form an \CrefAndHyperrefIfExist{definition:adjoint_functors_between_categories_unit_counit_of_adjoint_functors}{adjunction} $F \dashv G$, with \CrefAndHyperrefIfExist{definition:adjoint_functors_between_categories_unit_counit_of_adjoint_functors}{unit} $\eta : \mathrm{Id}_{\mathcal{C}} \implies GF$ and \CrefAndHyperrefIfExist{definition:adjoint_functors_between_categories_unit_counit_of_adjoint_functors}{counit} $\varepsilon : FG \implies \mathrm{Id}_{\mathcal{D}}$. Let $\mathcal{J}$ be a category.

    \begin{enumerate}
        \item \label{item:left_adjoint_preserves_colimits} Let $D : \mathcal{J} \to \mathcal{C}$ be a \CrefAndHyperrefIfExist{definition:diagram_in_a_category_indexed_by_a_category}{diagram} such that the \CrefAndHyperrefIfExist{definition:limit_and_colimit_of_a_diagram_in_a_category}{colimit} $\colim_{\mathcal{J}} D$ exists in $\mathcal{C}$, with universal cocone $(\lambda_j : D(j) \to \colim_{\mathcal{J}} D)_{j \in \Ob(\mathcal{J})}$. Then the family
        $$
        (F(\lambda_j) : F(D(j)) \to F(\colim_{\mathcal{J}} D))_{j \in \Ob(\mathcal{J})}
        $$
        is a \CrefAndHyperrefIfExist{definition:limit_and_colimit_of_a_diagram_in_a_category}{colimit cocone} for the diagram $F \circ D : \mathcal{J} \to \mathcal{D}$. In other words, for every object $Y \in \mathcal{D}$ and every family of morphisms $(\nu_j : F(D(j)) \to Y)_{j \in \Ob(\mathcal{J})}$ satisfying $\nu_k \circ F(D(\alpha)) = \nu_j$ for all morphisms $\alpha : j \to k$ in $\mathcal{J}$, there exists a unique morphism $f : F(\colim_{\mathcal{J}} D) \to Y$ such that $f \circ F(\lambda_j) = \nu_j$ for all $j \in \Ob(\mathcal{J})$.

        \item \label{item:right_adjoint_preserves_limits} Let $E : \mathcal{J} \to \mathcal{D}$ be a diagram such that the \CrefAndHyperrefIfExist{definition:limit_and_colimit_of_a_diagram_in_a_category}{limit} $\lim_{\mathcal{J}} E$ exists in $\mathcal{D}$, with universal cone $(\rho_j : \lim_{\mathcal{J}} E \to E(j))_{j \in \Ob(\mathcal{J})}$. Then the family
        $$
        (G(\rho_j) : G(\lim_{\mathcal{J}} E) \to G(E(j)))_{j \in \Ob(\mathcal{J})}
        $$
        is a \CrefAndHyperrefIfExist{definition:limit_and_colimit_of_a_diagram_in_a_category}{limit cone} for the diagram $G \circ E : \mathcal{J} \to \mathcal{C}$. In other words, for every object $X \in \mathcal{C}$ and every family of morphisms $(\sigma_j : X \to G(E(j)))_{j \in \Ob(\mathcal{J})}$ satisfying $G(E(\alpha)) \circ \sigma_j = \sigma_k$ for all morphisms $\alpha : j \to k$ in $\mathcal{J}$, there exists a unique morphism $g : X \to G(\lim_{\mathcal{J}} E)$ such that $G(\rho_j) \circ g = \sigma_j$ for all $j \in \Ob(\mathcal{J})$.
    \end{enumerate}
\end{theorem}

\begin{proof}
    \TODO{AI generated, verify}

    We prove part~\ref{item:left_adjoint_preserves_colimits}; part~\ref{item:right_adjoint_preserves_limits} follows by duality in the opposite categories $\mathcal{C}^{\op}$ and $\mathcal{D}^{\op}$, where $F$ and $G$ are swapped and all arrows reversed.

    Fix a diagram $D : \mathcal{J} \to \mathcal{C}$, its colimit $\colim_{\mathcal{J}} D$, and the universal cocone $(\lambda_j)_{j \in \Ob(\mathcal{J})}$. Let $Y \in \mathcal{D}$ and let $(\nu_j : F(D(j)) \to Y)_{j \in \Ob(\mathcal{J})}$ be a cocone under $F \circ D$. We must construct a unique morphism $f : F(\colim_{\mathcal{J}} D) \to Y$ with $f \circ F(\lambda_j) = \nu_j$ for all $j$:
    \[
    \begin{tikzcd}[row sep=large, column sep=large]
        F(D(j)) \arrow[rd, "F(\lambda_j)"'] \arrow[ddr, "\nu_j"', bend right=15] & \\
        & F(\colim_{\mathcal{J}} D) \arrow[d, dashed, "\exists ! f"] \\
        & Y
    \end{tikzcd}
    \]

    We first pull the cocone $(\nu_j)$ from $\mathcal{D}$ back to $\mathcal{C}$ using $G$ and the unit $\eta$. For each $j \in \Ob(\mathcal{J})$, define
    $$
    \mu'_j := G(\nu_j) \circ \eta_{D(j)} : D(j) \longrightarrow G(Y).
    $$
    We verify that $(\mu'_j)_j$ is a cocone from $D$ to $G(Y)$. Let $\alpha : j \to k$ be any morphism in $\mathcal{J}$. Since $(\nu_j)_j$ is a cocone under $F \circ D$, we have $\nu_k \circ F(D(\alpha)) = \nu_j$. We compute:
    $$
    \mu'_k \circ D(\alpha) = G(\nu_k) \circ \eta_{D(k)} \circ D(\alpha).
    $$
    By the naturality of $\eta$ applied to the morphism $D(\alpha) : D(j) \to D(k)$, we have $\eta_{D(k)} \circ D(\alpha) = GF(D(\alpha)) \circ \eta_{D(j)}$. Substituting:
    $$
    G(\nu_k) \circ GF(D(\alpha)) \circ \eta_{D(j)} = G(\nu_k \circ F(D(\alpha))) \circ \eta_{D(j)}.
    $$
    Using the cocone condition $\nu_k \circ F(D(\alpha)) = \nu_j$, this becomes $G(\nu_j) \circ \eta_{D(j)} = \mu'_j$. Thus $(\mu'_j)_j$ is a cocone.

    Next we invoke the universal property of $\colim_{\mathcal{J}} D$. Since $(\lambda_j)_j$ is a colimit cocone for $D$ and $(\mu'_j)_j$ is a cocone from $D$ to $G(Y)$, there exists a unique morphism
    $$
    u : \colim_{\mathcal{J}} D \longrightarrow G(Y)
    $$
    such that $u \circ \lambda_j = \mu'_j$ for all $j \in \Ob(\mathcal{J})$.

    We then push the morphism $u$ forward to $\mathcal{D}$ using the counit $\varepsilon$. Define
    $$
    f := \varepsilon_Y \circ F(u) : F(\colim_{\mathcal{J}} D) \longrightarrow Y,
    $$
    where $\varepsilon_Y : FG(Y) \to Y$ is the component of the counit at $Y$.

    We verify that $f$ mediates the cocone $(\nu_j)$. For each $j \in \Ob(\mathcal{J})$, we compute:
    $$
    f \circ F(\lambda_j) = \varepsilon_Y \circ F(u) \circ F(\lambda_j) = \varepsilon_Y \circ F(u \circ \lambda_j).
    $$
    Substituting $u \circ \lambda_j = \mu'_j = G(\nu_j) \circ \eta_{D(j)}$:
    $$
    \varepsilon_Y \circ F(G(\nu_j) \circ \eta_{D(j)}) = \varepsilon_Y \circ FG(\nu_j) \circ F(\eta_{D(j)}).
    $$
    By the naturality of $\varepsilon$ applied to the morphism $\nu_j : F(D(j)) \to Y$, we have the commutative square
    \[
    \begin{tikzcd}[row sep=large, column sep=large]
        FG(F(D(j))) \arrow[r, "FG(\nu_j)"] \arrow[d, "\varepsilon_{F(D(j))}"'] & FG(Y) \arrow[d, "\varepsilon_Y"] \\
        F(D(j)) \arrow[r, "\nu_j"'] & Y
    \end{tikzcd}
    \]
    That is, $\varepsilon_Y \circ FG(\nu_j) = \nu_j \circ \varepsilon_{F(D(j))}$. Substituting:
    $$
    f \circ F(\lambda_j) = \nu_j \circ \varepsilon_{F(D(j))} \circ F(\eta_{D(j)}).
    $$
    By the first triangle identity for the adjunction $F \dashv G$, applied to the object $D(j) \in \mathcal{C}$, we have $\varepsilon_{F(D(j))} \circ F(\eta_{D(j)}) = \mathrm{id}_{F(D(j))}$. Therefore:
    $$
    f \circ F(\lambda_j) = \nu_j \circ \mathrm{id}_{F(D(j))} = \nu_j.
    $$

    Finally we establish uniqueness. Suppose $f' : F(\colim_{\mathcal{J}} D) \to Y$ is another morphism satisfying $f' \circ F(\lambda_j) = \nu_j$ for all $j$. Define
    $$
    u' := G(f') \circ \eta_{\colim_{\mathcal{J}} D} : \colim_{\mathcal{J}} D \longrightarrow G(Y).
    $$
    For each $j$, we compute:
    $$
    u' \circ \lambda_j = G(f') \circ \eta_{\colim_{\mathcal{J}} D} \circ \lambda_j.
    $$
    By the naturality of $\eta$ applied to $\lambda_j : D(j) \to \colim_{\mathcal{J}} D$, we have $\eta_{\colim_{\mathcal{J}} D} \circ \lambda_j = GF(\lambda_j) \circ \eta_{D(j)}$. Hence:
    $$
    u' \circ \lambda_j = G(f') \circ GF(\lambda_j) \circ \eta_{D(j)} = G(f' \circ F(\lambda_j)) \circ \eta_{D(j)}.
    $$
    Since $f' \circ F(\lambda_j) = \nu_j$, this equals $G(\nu_j) \circ \eta_{D(j)} = \mu'_j$. Thus $u'$ is also a mediating morphism for the cocone $(\mu'_j)_j$ through $(\lambda_j)_j$. By the uniqueness of $u$, we have $u' = u$.

    It remains to show that $f' = f$. We compute:
    $$
    \varepsilon_Y \circ F(u') = \varepsilon_Y \circ F(G(f') \circ \eta_{\colim_{\mathcal{J}} D}) = \varepsilon_Y \circ FG(f') \circ F(\eta_{\colim_{\mathcal{J}} D}).
    $$
    By the naturality of $\varepsilon$ applied to $f' : F(\colim_{\mathcal{J}} D) \to Y$, we have $\varepsilon_Y \circ FG(f') = f' \circ \varepsilon_{F(\colim_{\mathcal{J}} D)}$. Hence:
    $$
    \varepsilon_Y \circ F(u') = f' \circ \varepsilon_{F(\colim_{\mathcal{J}} D)} \circ F(\eta_{\colim_{\mathcal{J}} D}).
    $$
    By the first triangle identity applied to $\colim_{\mathcal{J}} D$, we have $\varepsilon_{F(\colim_{\mathcal{J}} D)} \circ F(\eta_{\colim_{\mathcal{J}} D}) = \mathrm{id}_{F(\colim_{\mathcal{J}} D)}$. Therefore:
    $$
    \varepsilon_Y \circ F(u') = f' \circ \mathrm{id}_{F(\colim_{\mathcal{J}} D)} = f'.
    $$
    Since $u' = u$, we obtain $f' = \varepsilon_Y \circ F(u) = f$. This establishes uniqueness.

    The cocone $(F(\lambda_j))_j$ therefore satisfies the universal property of a colimit cocone for $F \circ D$, completing the proof of part~\ref{item:left_adjoint_preserves_colimits}.
\end{proof}
